\section{Contexto del evento}
Las Jornadas de Ciencia de Datos son un congreso de dos d\'ias, con cuatro franjas cada d\'ia (ocho en total), en tres espacios: Auditorio (200 plazas), Aula de ordenadores (40) y Sala de seminarios (80).

La organizaci\'on tiene 18 actividades candidatas, de tres categor\'ias: Machine Learning (ML), Estad\'istica e Ingenier\'ia de datos. Cuestan 12.050 en total, pero el presupuesto es de 6.500, as\'i que no se pueden programar todas. Hay que decidir qu\'e actividades se hacen y en qu\'e espacio y franja empiezan.

El contexto aporta de forma natural las condiciones del problema:
\begin{itemize}[nosep]
\item La conferencia inaugural y la de clausura son obligatorias.
\item Los talleres y el hackathon duran varias franjas y solo caben en el Aula de ordenadores.
\item Las charlas con m\'as p\'ublico no caben en la Sala de seminarios.
\item Los ponentes invitados solo est\'an disponibles en ciertas franjas.
\item Hay actividades que compiten por el mismo p\'ublico y no deben coincidir.
\end{itemize}

\section{Datos}
Los datos est\'an en 7 archivos CSV, separados del c\'odigo, para poder ejecutar otras instancias. Los valores (costes, asistencias, valoraciones) son estimaciones del grupo para este evento. Las tablas completas de la instancia est\'an en el Anexo A.

\begin{center}\small
\begin{tabular}{@{}ll@{}}\toprule
Archivo & Contenido\\\midrule
actividades.csv & Duraci\'on, categor\'ia, asistencia, valoraci\'on, coste y si es obligatoria\\
espacios.csv & Capacidad de cada espacio\\
compat\_espacios.csv & Espacios compatibles con cada actividad\\
disp\_actividades.csv & Franjas en las que puede celebrarse cada actividad (sin fila = no disponible)\\
disp\_espacios.csv & Franjas en las que est\'a abierto cada espacio\\
incompatibilidades.csv & Pares de actividades que no pueden coincidir en ninguna franja\\
parametros.csv & Presupuesto (6.500), pesos, n\'umero de franjas (8) y franjas por d\'ia (4)\\\bottomrule
\end{tabular}\end{center}

Las franjas van de 1 a 8: el d\'ia 1 son las franjas 1--4 y el d\'ia 2 las 5--8.

El problema tiene decisi\'on real porque no se puede programar todo sin incumplir alguna condici\'on: el coste total (12.050) supera el presupuesto (6.500) y las duraciones suman 27 franjas, m\'as de las 21 que ofrecen los espacios disponibles.

\section{Funci\'on objetivo y justificaci\'on}
Se maximiza una suma ponderada de tres criterios: asistencia total, valoraci\'on total y n\'umero de actividades programadas. Cada criterio se normaliza dividi\'endolo por su suma sobre todas las actividades, de modo que vale entre 0 y 1.
\[
\max \sum_{i \in I} y_i \left( w_A \frac{a_i}{A_{tot}} + w_V \frac{v_i}{V_{tot}} + w_N \frac{1}{N} \right)
\]
Aqu\'i $y_i$ vale 1 si la actividad $i$ se programa, $a_i$ es su asistencia esperada, $v_i$ su valoraci\'on y $N$ el n\'umero total de actividades. $A_{tot}$ y $V_{tot}$ son la suma de asistencias y de valoraciones de todas las actividades. Los pesos $w_A$, $w_V$ y $w_N$ suman 1.

\paragraph{Por qu\'e esta combinaci\'on.}
\begin{itemize}[nosep]
\item La asistencia mide el alcance del evento y la valoraci\'on mide su calidad. Un congreso busca las dos cosas, y suelen tirar en direcciones distintas: los talleres tienen mucha valoraci\'on y poco p\'ublico, y las charlas generales al rev\'es.
\item El n\'umero de actividades premia programar m\'as actividades con el mismo presupuesto y evita elegir pocas actividades muy caras.
\end{itemize}

\paragraph{Por qu\'e se normaliza.} La asistencia est\'a en cientos de personas y la valoraci\'on entre 1 y 5. Sin normalizar, la asistencia dominar\'ia y los pesos no ser\'ian comparables. Con la normalizaci\'on, un peso de 0,4 significa lo mismo en cualquier criterio y en cualquier instancia.

\paragraph{Pesos y escenarios.} Los pesos son par\'ametros modificables desde la aplicaci\'on. Los valores por defecto son 0,4 / 0,4 / 0,2. Los escenarios del trabajo se diferencian solo por estos pesos, priorizando la asistencia (1 / 0 / 0), la valoraci\'on (0 / 1 / 0) o el n\'umero de actividades (0,1 / 0,1 / 0,8). Con el presupuesto de 6.500 salen tres programaciones distintas: 9 actividades con asistencia total 880 y valoraci\'on total 30,5 (asistencia), 12 con 788 y 43,7 (valoraci\'on) y 12 con 863 y 42,3 (n\'umero de actividades). El n\'umero de actividades discrimina menos que los otros dos criterios, porque programar m\'as actividades tambi\'en suma asistencia y valoraci\'on.

\section{Ampliaci\'on y justificaci\'on}
Se a\~nade al modelo b\'asico una caracter\'istica: las incompatibilidades entre actividades (restricci\'on R9).

\textbf{Incompatibilidades entre actividades.} Dos actividades incompatibles pueden programarse las dos, pero nunca en la misma franja, ni siquiera en espacios distintos. Tiene sentido en un congreso porque hay actividades que comparten p\'ublico o ponente: un asistente no puede estar en las dos a la vez. Los cinco pares definidos son (3,10), (12,16), (9,17), (7,18) y (8,14).

La restricci\'on cambia la planificaci\'on y no solo la selecci\'on. Cuando dos actividades incompatibles tienen disponibilidad en una sola franja com\'un, como los ponentes invitados 7 y 18, que solo est\'an en la franja 7, el modelo no puede programar las dos. En las pruebas con el presupuesto de 6.500, quitar esta restricci\'on mejora el valor de la funci\'on objetivo en los cuatro escenarios probados, as\'i que condiciona de verdad el resultado.

\section{Formulaci\'on del modelo}
\paragraph{Conjuntos.}
\begin{itemize}[nosep]
\item $I$: actividades (1 a 18). $E$: espacios. $T=\{1,\dots,F\}$: franjas.
\item $O\subset I$: actividades obligatorias.
\item $P$: pares $(i,j)$ de actividades incompatibles.
\end{itemize}
\paragraph{Par\'ametros.}
\begin{itemize}[nosep]
\item $d_i$ duraci\'on (franjas), $a_i$ asistencia, $v_i$ valoraci\'on, $c_i$ coste de la actividad $i$.
\item $cap_e$ capacidad del espacio $e$. $B$ presupuesto.
\item $comp(i,e)=1$ si $i$ es compatible con $e$. $dispA(i,s)=1$ si $i$ puede celebrarse en la franja $s$. $dispE(e,s)=1$ si $e$ est\'a abierto en la franja $s$.
\item $F$ n\'umero de franjas (8) y $FD$ franjas por d\'ia (4), ambos en parametros.csv. $dia(t)=\lceil t/FD\rceil$.
\item $w_A,w_V,w_N$ pesos, con $w_A+w_V+w_N=1$.
\item $N$ n\'umero de actividades. $A_{tot}$ y $V_{tot}$: suma de asistencias y de valoraciones de todas las actividades.
\end{itemize}
\paragraph{Variables (binarias).}
\begin{itemize}[nosep]
\item $x_{i,e,t}=1$ si la actividad $i$ empieza en el espacio $e$ en la franja $t$.
\item $y_i=1$ si la actividad $i$ se programa.
\end{itemize}
\paragraph{Funci\'on objetivo.}
\[
\max \sum_{i \in I} y_i \left( w_A \frac{a_i}{A_{tot}} + w_V \frac{v_i}{V_{tot}} + w_N \frac{1}{N} \right)
\]
\paragraph{Restricciones.}

\textbf{R1.} Cada actividad se programa como m\'aximo una vez, y $y_i$ indica si se programa:
\[ y_i=\sum_{e\in E}\sum_{t\in T}x_{i,e,t}\le 1\qquad\forall i\in I \]
\textbf{R2.} Las obligatorias se programan siempre:
\[ y_i=1\qquad\forall i\in O \]
\textbf{R3.} Solo en espacios compatibles:
\[ x_{i,e,t}=0\quad\text{si } comp(i,e)=0 \]
\textbf{R4.} La asistencia cabe en el espacio:
\[ x_{i,e,t}=0\quad\text{si } a_i>cap_e \]
\textbf{R5.} Actividad y espacio disponibles en todas las franjas que ocupa, no solo en la de inicio:
\[ x_{i,e,t}\le dispA(i,s)\cdot dispE(e,s)\qquad\forall s\in\{t,\dots,t+d_i-1\} \]
\textbf{R6.} No terminar despu\'es de la \'ultima franja ni cruzar de un d\'ia a otro:
\[ x_{i,e,t}=0\quad\text{si } t+d_i-1>F \ \text{ o } \ dia(t)\neq dia(t+d_i-1) \]
\textbf{R7.} Presupuesto:
\[ \sum_{i\in I}c_i\,y_i\le B \]
\textbf{R8.} Sin solapes en un espacio (una actividad que empieza en $t$ ocupa las franjas $t,\dots,t+d_i-1$):
\[ \sum_{i\in I}\ \sum_{t=\max(1,\,s-d_i+1)}^{s}x_{i,e,t}\le 1\qquad\forall e\in E,\ \forall s\in T \]
\textbf{R9 (ampliaci\'on).} Dos actividades incompatibles no coinciden en ninguna franja, aunque est\'en en espacios distintos:
\[ \sum_{e\in E}\sum_{t=\max(1,\,s-d_i+1)}^{s}x_{i,e,t}+\sum_{e\in E}\sum_{t=\max(1,\,s-d_j+1)}^{s}x_{j,e,t}\le 1\qquad\forall (i,j)\in P,\ \forall s\in T \]

\clearpage
\appendix
\section{Datos de la instancia}
\textbf{Actividades} (E1 Auditorio, E2 Aula de ordenadores, E3 Sala de seminarios; franjas 1--4 d\'ia 1, 5--8 d\'ia 2)

{\footnotesize
\begin{longtable}{@{}r p{4.6cm} c l r r r c l l@{}}
\toprule
Id & Actividad & Dur. & Categor\'ia & Asist. & Valor. & Coste & Obl. & Esp. & Franjas\\\midrule\endhead
1 & Conferencia inaugural: IA y sociedad & 1 & ML & 180 & 4,2 & 1.200 & S\'i & E1 & 1, 2\\
2 & Conferencia de clausura: datos y \'etica & 1 & Estad\'istica & 160 & 4,0 & 1.000 & S\'i & E1 & 7, 8\\
3 & Taller de redes neuronales & 3 & ML & 35 & 4,8 & 900 & No & E2 & 1,2,3,5--8\\
4 & Taller de pandas y limpieza de datos & 2 & Ingenier\'ia & 38 & 4,2 & 500 & No & E2 & todas\\
5 & Hackathon de datos & 3 & Ingenier\'ia & 40 & 4,9 & 1.100 & No & E2 & 5--8\\
6 & Taller de inferencia bayesiana & 2 & Estad\'istica & 30 & 4,7 & 550 & No & E2 & todas\\
7 & Charla: IA generativa & 1 & ML & 190 & 3,6 & 1.800 & No & E1 & 7\\
8 & Charla: pipelines de datos en producci\'on & 1 & Ingenier\'ia & 45 & 3,0 & 350 & No & E1, E3 & todas\\
9 & Charla: series temporales & 1 & Estad\'istica & 40 & 3,0 & 250 & No & E1, E3 & 4\\
10 & Mesa redonda: datos y empleo & 2 & Ingenier\'ia & 110 & 3,4 & 900 & No & E1 & 2, 3, 4\\
11 & Charla: visualizaci\'on de datos & 1 & Estad\'istica & 35 & 2,8 & 250 & No & E1, E3 & todas\\
12 & Charla: aprendizaje por refuerzo & 1 & ML & 50 & 3,2 & 350 & No & E1, E3 & 6, 7, 8\\
13 & Taller de SQL avanzado & 2 & Ingenier\'ia & 30 & 4,0 & 400 & No & E2 & todas\\
14 & Charla: Big Data con Spark & 1 & Ingenier\'ia & 60 & 3,5 & 450 & No & E1, E3 & 1--4\\
15 & Charla: contrastes y A/B testing & 1 & Estad\'istica & 30 & 2,7 & 200 & No & E1, E3 & todas\\
16 & Charla: explicabilidad en ML & 1 & ML & 85 & 3,6 & 500 & No & E1 & 6, 7, 8\\
17 & Seminario de estad\'istica espacial & 2 & Estad\'istica & 45 & 3,0 & 350 & No & E1, E3 & 3, 4\\
18 & Charla: \'etica y sesgos algor\'itmicos & 1 & ML & 75 & 4,0 & 1.000 & No & E1, E3 & 7\\\bottomrule
\end{longtable}}

La columna Esp. indica donde puede celebrarse cada actividad una vez aplicadas la compatibilidad y la capacidad: las actividades 7, 10 y 16 son compatibles con la Sala de seminarios por tipo, pero su asistencia supera las 80 plazas.

\textbf{Espacios}

\begin{tabular}{@{}l l r l@{}}\toprule
Id & Espacio & Capacidad & Franjas abiertas\\\midrule
E1 & Auditorio & 200 & 1 a 8\\
E2 & Aula de ordenadores & 40 & 1, 2, 3, 5, 6, 7, 8 (cerrada en la 4)\\
E3 & Sala de seminarios & 80 & 2, 3, 4, 6, 7, 8 (cerrada en la 1 y la 5)\\\bottomrule
\end{tabular}

\textbf{Incompatibilidades:} (3,10), (12,16), (9,17), (7,18) y (8,14).

\textbf{Presupuesto:} 6.500. \textbf{Pesos por defecto:} $w_A=0{,}4$, $w_V=0{,}4$, $w_N=0{,}2$. \textbf{Franjas:} 8, con 4 por d\'ia.
